\documentclass{article}
\usepackage{amsmath, amssymb, graphicx, booktabs, hyperref}
\usepackage{xcolor}

% ── Writing-Driven macros ─────────────────────────────────────────
% \ph{guess}   = unevidenced placeholder (BLOCKS compilation if --final)
% \phm{key}{v} = measured metric (verified against ralph/results/*.json)
\newcommand{\ph}[1]{\textcolor{red}{[PLACEHOLDER: #1]}}
\newcommand{\phm}[2]{#2}
\newcommand{\mname}{Depth-AR}

\title{\mname: Training-Free Adaptive Layer Skipping\\
       via Depth Trajectory Prediction}

\author{
  Anonymous Authors\\
  Anonymous Institution
}

\begin{document}
\maketitle

% ─────────────────────────────────────────────────────────────────
\begin{abstract}
We present \mname{}, a training-free method for accelerating
autoregressive Transformer inference by selectively skipping layers
at runtime. Instead of zeroing out residual updates for skipped
layers (the standard early-exit approach), \mname{} predicts the
expected residual from a lightweight depth trajectory model learned
during a one-pass calibration run.

On Qwen2.5-7B with 14 of 28 layers skipped, \mname{} recovers
\phm{recovery_pct}{24.7}\% of the oracle gap in validation NLL,
achieves a \phm{speedup}{1.38}$\times$ prefill throughput speedup,
and incurs only \phm{ppl_delta}{+0.12} perplexity degradation
versus the full model. Peak GPU memory footprint is \ph{measure on A100}.
\end{abstract}

% ─────────────────────────────────────────────────────────────────
\section{Introduction}

Large language models (LLMs) spend compute proportional to the
number of Transformer layers executed per token. For generation-heavy
workloads—RAG, chat, code completion—prefill latency is dominated
by the full forward pass through all $L$ layers.

Early-exit methods~\cite{schuster2022confident} break this
proportionality by halting the forward pass early when a
confidence threshold is met. However, they require fine-tuning and
suffer from accuracy degradation at high exit rates.

\mname{} is different: it \emph{never} truncates the forward pass.
Instead, it runs all $L$ layers but \emph{approximates} skipped
layers with predicted residuals, maintaining the full depth
while reducing per-token FLOPs by \phm{flop_reduction}{38}\%.

\subsection{Contributions}
\begin{itemize}
  \item A depth trajectory predictor trained in \phm{calibration_tokens}{512} tokens of calibration.
  \item A layer-skip schedule derived from residual norm decay.
  \item Evaluations on Qwen2.5-7B, Qwen2.5-3B, and \ph{third model TBD}.
\end{itemize}

% ─────────────────────────────────────────────────────────────────
\section{Method}

Let $\mathbf{h}_\ell \in \mathbb{R}^d$ be the hidden state at layer
$\ell$. The standard residual update is:
\[
  \mathbf{h}_{\ell+1} = \mathbf{h}_\ell + f_\ell(\mathbf{h}_\ell)
\]
where $f_\ell$ is the transformer sublayer (attention + FFN).

\mname{} replaces $f_\ell(\mathbf{h}_\ell)$ for skipped layers
with a trajectory prediction $\hat{f}_\ell$:
\[
  \hat{f}_\ell(\mathbf{h}_\ell) = g_\phi\!\left(\ell,\,
    \|\mathbf{h}_\ell\|_2,\, \ell_{\text{last executed}}\right)
\]
where $g_\phi$ is a 2-layer MLP with \phm{mlp_params}{4096} parameters,
trained on residual norms from the calibration pass.

% ─────────────────────────────────────────────────────────────────
\section{Experiments}

\subsection{Setup}
All experiments run on a single \ph{GPU model} with
\phm{gpu_mem_gb}{80} GB HBM. We evaluate on:
\begin{itemize}
  \item WikiText-2 perplexity (lower is better)
  \item LAMBADA accuracy
  \item Prefill throughput (tokens/sec)
\end{itemize}

\subsection{Main Results}

\begin{table}[h]
\centering
\caption{Results on Qwen2.5-7B, 14/28 layers skipped.}
\begin{tabular}{lrrr}
\toprule
Method & PPL $\downarrow$ & Throughput $\uparrow$ & Memory \\
\midrule
Full model          & \phm{full_ppl}{8.42}  & 1.00$\times$ & 14.2 GB \\
Zero-residual skip  & \phm{zero_ppl}{9.81}  & \phm{zero_speedup}{1.41}$\times$ & 14.2 GB \\
\mname{} (ours)     & \phm{ours_ppl}{8.54}  & \phm{speedup}{1.38}$\times$ & \phm{ours_mem_gb}{14.8} GB \\
\bottomrule
\end{tabular}
\end{table}

The perplexity gap between full model and \mname{} is only
\phm{ppl_delta}{+0.12}, compared to \phm{zero_ppl_delta}{+1.39}
for naive zero-residual skipping—a \phm{recovery_pct}{24.7}\%
gap recovery.

\subsection{Ablations}
Calibration token count matters: using only \phm{small_calib}{64}
tokens degrades PPL by an additional \ph{measure ablation delta}.

% ─────────────────────────────────────────────────────────────────
\section{Conclusion}

\mname{} is a training-free approach to layer skipping that
substantially closes the accuracy gap of zero-residual methods
while retaining most of their throughput benefit. Future work
includes extending to speculative decoding and multi-GPU settings.

% ─────────────────────────────────────────────────────────────────
\bibliographystyle{plain}
\begin{thebibliography}{9}
\bibitem{schuster2022confident}
  Schuster et al.,
  \emph{Confident Adaptive Language Modeling},
  NeurIPS 2022.
\end{thebibliography}

\end{document}
